\section{Discussion}
\label{sec:discussion}

\hi{Why the CNN pair helps and mixed ensembles do not.}
A vote can only fix an error if enough members get that image right.
The two CNNs are both accurate, and although their $Q$ is high they still disagree on a fair number of images (57 and 51 one-sided errors on the balanced set), so averaging settles most disagreements in favour of the right class.
A hand-crafted member is diverse in the wrong way.
It is wrong on many images the CNNs already handle, so in a pair it can only pull the stronger model down, and in a larger vote it can outvote it.
This fits the familiar view that ensemble members need to be both accurate and diverse~\cite{dietterich2000ensemble,kuncheva2003diversity}; here the accuracy gap between the families outweighs whatever the weak members add.
A combiner learned on validation predictions might exploit the high oracle accuracy, but we have not tried one.

\hi{Why augmentation hurts the CNNs.}
To bring every class to 544 images, the rebalanced set generates up to 15 copies of each real tail image; \emph{River}, for example, has 34 real images and 510 synthetic ones.
The hand-crafted descriptors are histograms that barely change under rotation and mild photometric jitter, so for them the copies mostly correct the class prior.
For the CNNs, one explanation consistent with our results is that they fit the repeated content of the few real tail images.
Their tail recall drops below that of plain long-tailed training, whereas loss re-weighting, which shifts the prior without duplicating any content, improves both networks.
The CNNs also apply random flips and rotations during training in every setting, so the offline copies add less new variation for them than for the hand-crafted pipelines.
We did not try smaller augmentation budgets, and we evaluated the weighted loss for the CNNs only.

\hi{The tie-breaking pitfall.}
The earlier version of this study broke voting ties by member order.
Every two-member ensemble therefore matched its first member exactly, which made it look as though ensembles ``collapse to their weakest member''.
Nothing of the sort was happening; with two voters every disagreement is a tie, and the rule decided all of them the same way.
Any hard vote with an even number of members needs an explicit tie-break, and soft voting sidesteps the problem.

\hi{Limitations.}
All results come from one dataset with $256\times256$ images and 15 classes, a single (alphabetical) long-tail ordering, and one training run per configuration, so the bootstrap intervals do not capture training variance.
We fine-tuned \eb at $224\times224$ rather than its native $300\times300$ and unfroze only its last stage, which may understate it relative to the fully fine-tuned \rn.
The Grad-CAM analysis is qualitative and covers a single confusion.
